\cvsection{Experience}

\begin{cventries}
  \cventry
    {Python Engineer $\rightarrow$ ML Engineer $\rightarrow$ Senior ML Engineer}
    {Creative Fabrica}
    {Amsterdam, the Netherlands}
    {Dec. 2023 -- present}
    {
      \begin{cvitems}
      \item Own the generative AI \textbf{Python} codebase, infrastructure, and internal testing and QA tool.
      \item Developed an AI agent for canvas manipulation and built an evaluation pipeline to benchmark against the baseline agent.
      \item Developed and deployed to production an inference pipeline for content matching.
      \item Designed and maintained distributed event-driven architecture using \textbf{Kafka}, \textbf{Redis} queues, and \textbf{Kubernetes}.
      \item Redesigned the architecture for image generation backend on \textbf{AWS}, significantly improving scalability and system performance.
      \end{cvitems}
    }
  \cventry
    {Data Scientist (contract)}
    {SBDA Group}
    {Remote}
    {Sep. 2023 -- Oct 2023}
    {
      \begin{cvitems}
      \item Designed and implemented an algorithmic solution in Python for optimizing a manufacturing supply chain process.
      \end{cvitems}
    }
  \cventry
    {Research fellow}
    {Artificial Intelligence Research Institute (AIRI)}
    {Moscow, Russia}
    {Feb. 2022 -- Apr. 2023}
    {
      \begin{cvitems}
      \item Built and trained a deep learning model to predict speech from brain data in \textbf{PyTorch}.
      \item Developed real-time \textbf{Python} pipeline for EEG timeseries prediction using Kalman filtering.
      \item Developed a \textbf{Python} library for type-safe timeseries arrays manipulation.
      \end{cvitems}
    }
  \cventry
    {Junior research fellow}
    {Higher School of Economics, Centre for Cognition and Decision Making}
    {Moscow, Russia}
    {Feb. 2015 -- Dec. 2021}
    {
      \begin{cvitems}
      % \item{Built a \href{https://github.com/dmalt/PSIICOS}{toolbox} for brain connectivity estimation and visualization in \textbf{MATLAB}.}
      \item Accelerated various parts of my PhD codebase to enable Monte-Carlo statistics, using \textbf{CUDA} (cuBLAS) and \textbf{C}.
      \item{Led 4-people group on a project for real-time ML feature extraction and GPU-accelerated 3D visualization of brain activity in \textbf{Python}.}
            % Technologies used: \textbf{Python}, \textbf{PyQt}, \textbf{MNE-python}, \textbf{Vispy}, \textbf{Numpy}.}
      \item{Sped up by a factor of 10 the beamformer inverse solver in \href{https://github.com/mne-tools/mne-python}{MNE-python} by modifying the algorithm for vectorized computations.}
      % \item{Created a pipeline for semi-automated MEG data preprocessing, see \href{https://github.com/dmalt/metacognition}{metacognition}, \href{https://github.com/dmalt/MRI_metacognition}{MRI\_metacognition}} on GitHub
      \end{cvitems}
    }
  % \cventry
  %   {Senior lecturer}
  %   {Higher School of Economics, Faculty of Computer Science}
  %   {Moscow, Russia}
  %   {Aug. 2019 -- Dec. 2021}
  %   {
  %     \begin{cvitems}
  %     \item Developed and taught a course ``\textbf{MATLAB} for data analysis'' to undergraduate students
  %     \end{cvitems}
  %   }
  \cventry
    {Visiting researcher, temporary position}
    {University of Montreal, CERNEC lab.}
    {Montreal, Canada}
    {Oct. -- Dec. 2015, May 2016 -- Dec. 2016}
    {
      \begin{cvitems}
          \item{Co-developed an open-source \textbf{Python} package for GPU-accelerated brain data visualization \href{https://github.com/EtienneCmb/visbrain}{Visbrain}, \cite{Combrisson}.}
      \end{cvitems}
    }
\end{cventries}
